% ============================================================================
%  All tables for the paper.
%  Required packages: booktabs, xcolor, colortbl, multirow, siunitx (optional)
%
%  Color macros (define once in main preamble):
%      \definecolor{bestcolor}{HTML}{C62828}    % deep red for best
%      \definecolor{secondcolor}{HTML}{1565C0}  % deep blue for second-best
%      \definecolor{failcolor}{HTML}{FFCDD2}    % light red row tint for KL>0.1
%      \definecolor{passcolor}{HTML}{C8E6C9}    % light green row tint for KL<0.1
%      \newcommand{\best}[1]{\textcolor{bestcolor}{\textbf{#1}}}
%      \newcommand{\second}[1]{\textcolor{secondcolor}{\underline{#1}}}
% ============================================================================

% ----------------------------------------------------------------------------
% TABLE 1 — Master single-row-per-model summary (the decisive table).
% Place at the opening of §4 Experiments.
% ----------------------------------------------------------------------------
\begin{table*}[t]
\centering
\caption{%
\textbf{Cross-model summary.}
$l^{*}$: target hidden-state layer from Exp~1. $\alpha = \|\theta_{\text{DIM}}\|_{2}$:
intervention magnitude. ASR$_{\text{abl}}$: in-distribution Answer Switching Rate
under ablation. ``Surgical'' columns mark whether $\mathrm{mean\,KL} < 0.1$ on the
retain set. Higher ASR ($\uparrow$) and lower KL ($\downarrow$) are better.
The three models that fail \emph{both} axes at $k=1$ (\colorbox{failcolor!60}{red row tint})
all recover \emph{both} at $k=2$.}
\label{tab:master}
\small
\begin{tabular}{lrrrrrrcc}
\toprule
Model & $l^{*}$ & $\alpha$
      & \multicolumn{1}{c}{DIM ASR$_{\text{abl}}\!\uparrow$}
      & \multicolumn{1}{c}{TDO ASR$_{\text{abl}}\!\uparrow$}
      & \multicolumn{1}{c}{$k{=}2$ ASR$\!\uparrow$}
      & \multicolumn{1}{c}{$k{=}2$ KL$\!\downarrow$}
      & Surg.\ $k{=}1$ & Surg.\ $k{=}2$ \\
\midrule
Qwen-2.5-1.5B-Instruct & 17 &  13.05 & 0.010 & \best{1.000} & 1.000 & 0.008 & \checkmark & \checkmark \\
Gemma-2-2B-it          & 15 &  84.52 & 0.216 & \best{1.000} & 1.000 & 0.004 & \checkmark & \checkmark \\
\rowcolor{failcolor!50}
Qwen-2.5-7B-Instruct   & 18 &  26.61 & 0.000 & 0.020        & \best{1.000} & 0.037 & $\times$ & \checkmark \\
Gemma-2-9B-it          & 24 & 106.56 & 0.000 & 0.774        & \best{1.000} & 0.001 & \checkmark & \checkmark \\
\rowcolor{failcolor!50}
Llama-3.1-8B-Instruct  & 14 &   4.07 & 0.020 & \best{1.000} & 1.000 & 0.017 & $\times$ & \checkmark \\
\rowcolor{failcolor!50}
Qwen-2.5-14B-Instruct  & 30 &  48.25 & 0.000 & 0.658        & \best{1.000} & 0.009 & $\times$ & \checkmark \\
\bottomrule
\end{tabular}
\end{table*}


% ----------------------------------------------------------------------------
% TABLE 2 — Exp 2: TDO vs DIM at d=1, in-dist and OOD
% Place in §4.2.
% ----------------------------------------------------------------------------
\begin{table*}[t]
\centering
\caption{%
\textbf{Exp~2 — Truth Direction Optimization (TDO) versus Difference-in-Means
(DIM) at $k=1$.} ID: in-distribution (train/val split per source). OOD: train
on two of \{cities, animals, elements\}, evaluate on the held-out third.
\best{Best} per (model, intervention, split). The single TDO failure case is
Qwen-2.5-7B (italicized): at $k=1$, neither method achieves meaningful ablation
ASR; this motivates §\ref{sec:tco}.}
\label{tab:tdo_vs_dim}
\small
\begin{tabular}{l rrrr rrrr}
\toprule
& \multicolumn{4}{c}{In-distribution (ID)} & \multicolumn{4}{c}{Out-of-distribution (OOD)} \\
\cmidrule(lr){2-5} \cmidrule(lr){6-9}
Model
 & DIM abl & TDO abl & DIM add & TDO add
 & DIM abl & TDO abl & DIM add & TDO add \\
\midrule
Qwen-2.5-1.5B   & 0.010 & \best{1.000} & 1.000 & 1.000          & 0.133 & \best{0.998} & 1.000 & 1.000 \\
Gemma-2-2B      & 0.216 & \best{1.000} & 1.000 & 1.000          & 0.071 & \best{1.000} & 1.000 & 1.000 \\
\textit{Qwen-2.5-7B} & 0.000 & 0.020   & \best{0.995} & 0.940   & 0.015 & \best{0.080} & \best{0.974} & 0.813 \\
Gemma-2-9B      & 0.000 & \best{0.774} & 0.714 & \best{0.960}   & 0.045 & \best{0.327} & 0.837 & \best{0.963} \\
Llama-3.1-8B    & 0.020 & \best{1.000} & 1.000 & 1.000          & 0.185 & \best{1.000} & 0.985 & \best{1.000} \\
Qwen-2.5-14B    & 0.000 & \best{0.658} & 1.000 & 1.000          & 0.011 & \best{0.125} & 1.000 & 1.000 \\
\bottomrule
\end{tabular}
\end{table*}


% ----------------------------------------------------------------------------
% TABLE 3 — Exp 3: Cone ASR per (model, k), ID and OOD.
% Place in §4.3.
% ----------------------------------------------------------------------------
\begin{table*}[t]
\centering
\caption{%
\textbf{Exp~3 — Mean MC ASR (ablation) over $256$ samples drawn from the
learned cone.}
For all six models, $k=2$ suffices to reach $\mathrm{ASR}=1.0$
in-distribution. The standard deviation column (ID only, $k=2$) is $0$ for all
models, indicating uniform mediation across the cone interior (Claim~C2').
$k^{*}$ is the smallest $k$ at which the model is simultaneously surgical
(Tab.~\ref{tab:kl}) and saturated on ASR.}
\label{tab:cone_asr}
\small
\begin{tabular}{l rrrr rrrr c}
\toprule
& \multicolumn{4}{c}{In-distribution (ID)} & \multicolumn{4}{c}{Out-of-distribution (OOD)} & \\
\cmidrule(lr){2-5} \cmidrule(lr){6-9}
Model & $k{=}1$ & $k{=}2$ & $k{=}3$ & $k{=}4$ & $k{=}1$ & $k{=}2$ & $k{=}3$ & $k{=}4$ & $k^{*}$ \\
\midrule
Qwen-2.5-1.5B & 1.000 & 1.000 & 1.000 & 1.000 & 0.998 & 0.998 & 0.999 & 0.987 & 2 \\
Gemma-2-2B    & 1.000 & 1.000 & 1.000 & 1.000 & 1.000 & 1.000 & 0.997 & 0.939 & 2 \\
Qwen-2.5-7B   & 0.020 & \best{1.000} & 1.000 & 1.000 & 0.080 & 0.682 & 0.987 & 0.993 & 2 \\
Gemma-2-9B    & 0.774 & \best{1.000} & 1.000 & 1.000 & 0.327 & 1.000 & 0.980 & 0.946 & 2 \\
Llama-3.1-8B  & 1.000 & 1.000 & 0.807 & 0.983 & 1.000 & 1.000 & 0.916 & 1.000 & 2 \\
Qwen-2.5-14B  & 0.658 & \best{1.000} & 1.000 & 0.994 & 0.125 & 0.998 & 0.938 & 0.894 & 2 \\
\bottomrule
\end{tabular}
\end{table*}


% ----------------------------------------------------------------------------
% TABLE 4 — Exp 4: DIM-vs-cone cosine alignment at reference k=4
% Place in §4.4.
% ----------------------------------------------------------------------------
\begin{table}[t]
\centering
\caption{%
\textbf{Exp~4 — $|\cos(\theta_{\text{DIM}}, v_i)|$ at reference $k=4$.}
The first basis vector at $k{=}1$ remains near-aligned with DIM (right column,
warm-start sanity check), but in the full 4-dimensional cone every basis vector
$v_i$ has $|\cos|<0.1$ for every model: the cone captures \emph{novel}
structure orthogonal to the linear probe.}
\label{tab:orthogonality}
\small
\setlength{\tabcolsep}{4pt}
\begin{tabular}{l rrrr c r}
\toprule
& \multicolumn{4}{c}{$|\cos(\theta_{\text{DIM}}, v_i)|$ at $k=4$} & & $|\cos|$ at \\
Model & $v_1$ & $v_2$ & $v_3$ & $v_4$ & Interp.\ & $k=1$ \\
\midrule
Qwen-2.5-1.5B & 0.006 & 0.030 & 0.019 & 0.050 & no align       & 0.932 \\
Gemma-2-2B    & 0.031 & 0.096 & 0.012 & 0.070 & no align       & 0.905 \\
Qwen-2.5-7B   & 0.012 & 0.063 & 0.061 & 0.011 & no align       & 0.868 \\
Gemma-2-9B    & 0.084 & 0.017 & 0.017 & 0.009 & $v_1$ only     & 0.860 \\
Llama-3.1-8B  & 0.088 & 0.037 & 0.010 & 0.041 & $v_1$ only     & 0.842 \\
Qwen-2.5-14B  & 0.060 & 0.005 & 0.044 & 0.021 & $v_1$ only     & 0.816 \\
\bottomrule
\end{tabular}
\end{table}


% ----------------------------------------------------------------------------
% TABLE 5 — Exp 5: Mean KL on retain set per (model, k)
% Place in §4.5.
% ----------------------------------------------------------------------------
\begin{table}[t]
\centering
\caption{%
\textbf{Exp~5 — Mean Monte-Carlo KL divergence on the retain set.}
$32$ directions per cone $\times$ $200$ Alpaca prompts $\times$ last $30$ token
positions. Bold + light-red highlighting marks non-surgical entries
($\mathrm{KL} \ge 0.1$, Arditi et al.~2024). The same three models that fail
ASR at $k{=}1$ (Tab.~\ref{tab:tdo_vs_dim}) also fail surgicality there, and
all six pass at $k{=}2$.}
\label{tab:kl}
\small
\setlength{\tabcolsep}{6pt}
\begin{tabular}{l rrrr c}
\toprule
Model & $k{=}1$ & $k{=}2$ & $k{=}3$ & $k{=}4$ & All surg.? \\
\midrule
Qwen-2.5-1.5B & \cellcolor{passcolor!50}0.042 & \cellcolor{passcolor!50}0.008 & \cellcolor{passcolor!50}0.009 & \cellcolor{passcolor!50}0.008 & \checkmark \\
Gemma-2-2B    & \cellcolor{passcolor!50}0.028 & \cellcolor{passcolor!50}0.004 & \cellcolor{passcolor!50}0.007 & \cellcolor{passcolor!50}0.012 & \checkmark \\
Gemma-2-9B    & \cellcolor{passcolor!50}0.048 & \cellcolor{passcolor!50}0.001 & \cellcolor{passcolor!50}0.001 & \cellcolor{passcolor!50}0.002 & \checkmark \\
Llama-3.1-8B  & \cellcolor{failcolor!60}\textbf{0.282} & \cellcolor{passcolor!50}0.017 & \cellcolor{failcolor!60}\textbf{0.811} & \cellcolor{passcolor!50}0.017 & $\times$ \\
Qwen-2.5-14B  & \cellcolor{failcolor!60}\textbf{0.194} & \cellcolor{passcolor!50}0.009 & \cellcolor{passcolor!50}0.027 & \cellcolor{passcolor!50}0.010 & $\times$ \\
Qwen-2.5-7B   & \cellcolor{failcolor!60}\textbf{0.292} & \cellcolor{passcolor!50}0.037 & \cellcolor{failcolor!60}\textbf{0.202} & \cellcolor{passcolor!50}0.044 & $\times$ \\
\bottomrule
\end{tabular}
\end{table}


% ----------------------------------------------------------------------------
% TABLE 6 — Hyperparameters (consider moving to supplementary)
% ----------------------------------------------------------------------------
\begin{table}[t]
\centering
\caption{%
\textbf{Hyperparameters for TDO and TCO.} Inherited from
Wollschläger et al.~(2025) Table~3 and held fixed across all models.}
\label{tab:hyperparams}
\small
\begin{tabular}{ll}
\toprule
Component & Value \\
\midrule
Effective batch size            & 16 (gradient accumulation, physical batch 1) \\
Base learning rate              & $10^{-2}$ \\
LR schedule                     & plateau-reduce $\times 0.1$, patience 5, max 2 reductions \\
Optimizer                       & AdamW, weight decay 0 \\
Gradient clipping               & max-norm 1.0; spherical projection (TDO only) \\
$\lambda_{\text{abl}}$          & 1.0 \\
$\lambda_{\text{add}}$          & 0.2 \\
$\lambda_{\text{ret}}$          & 1.0 \\
$\alpha$ (intervention magnitude) & $\|\theta_{\text{DIM}}\|_{2}$ per model \\
MC samples / accumulation step  & 16 \\
Effective MC samples / batch    & 256 \\
Cone re-orthogonalization       & modified Gram--Schmidt every step \\
KL surgicality threshold        & 0.1 \\
Retain set size                 & 200 Alpaca prompts (factual-keyword filtered) \\
KL evaluation positions         & last 30 tokens per prompt \\
\bottomrule
\end{tabular}
\end{table}
